\documentclass[10pt, titlepage]{article}

\usepackage[margin=1in]{geometry}
\usepackage{fancyhdr}
\usepackage{amsfonts}
\usepackage{amsmath}
\usepackage{amssymb}
\usepackage{booktabs}
%\usepackage{tikz}
\usepackage{float}
\usepackage{bm}
\usepackage{pgfplots}


\title{\normalsize PLUS Session Worksheet \\[1em] \Large \textit{MATH 1554:} Linear Algebra}
\author{\textit{Covering Section:} 1.7}
\date{Fall 2023}

% ------- Set up header -------
\renewcommand{\sectionmark}[1]{\markright{#1}} % from https://tex.stackexchange.com/a/317393/227665
\pagestyle{fancy}
\fancyhead[L]{A. M., Fall 2023}
%\fancyhead[C]{\textit{\nouppercase{\rightmark}}}
\fancyhead[C]{\textit{Linear Independence (1.7)}}
\fancyhead[R]{\textbf{\thepage}}
% -------------------------------------------------------

% ------- Set up question environment -------
\newcounter{question}
\setcounter{question}{1}
\newenvironment{newquestion}
{
	\noindent\textbf{Question \thequestion}
}
{
\stepcounter{question}
\par
}
% -------------------------------------------------------

% ------- Set up Main Idea box environment -------
\newcounter{mainideas}
\setcounter{mainideas}{1}
\newcommand{\mainidea}{\textbf{Main Idea \#\themainideas \stepcounter{mainideas}}}
\setlength\fboxsep{1em}
%\newcommand{\midea}[2]{\begin{center}\fbox{\parbox{#1}{\centering \mainidea #2}} \end{center}} % Could be environment but I couldn't bother

% The way this works is that the savebox collects whatever was between the minipage and the end of it, then in the last line of the definition of the environment
% we regurgigate whatever was collected into an fbox which we center.
\newsavebox{\mybox}
\newenvironment{mideea}[1]
{%
\begin{lrbox}{\mybox}%
	\begin{minipage}{#1\textwidth}
		{\centering \mainidea \par}

}
{\end{minipage}\end{lrbox}\begin{center}\fbox{\usebox{\mybox}}\end{center}}
% -------------------------------------------------------


% ------- Set up math stuff -------
\renewcommand{\vec}[1]{\bm{#1}}
\DeclareMathOperator{\vspan}{Span}
\DeclareMathOperator{\row}{row}
\DeclareMathOperator{\col}{col}
\DeclareMathOperator{\colspace}{Col}
\DeclareMathOperator{\rowspace}{Row}
\DeclareMathOperator{\nullspace}{Nul}
\DeclareMathOperator{\rank}{rank}
\DeclareMathOperator{\proj}{proj}
% -------------------------------------------------------

% ------- Set up circle -------
\newcommand{\bubble}{\tikz[baseline=-0.625ex]\draw[radius=0.5em, fill=white] (0,0) circle;}
\newcommand{\dbubble}{\tikz[baseline=-0.625ex]\draw[radius=0.5em, fill=black] (0,0) circle;}
% -------------------------------------------------------

\makeatletter
\setlength{\@fptop}{0pt}
\makeatother

\begin{document}
\maketitle

\section*{\centering Linear Independence (1.7)}
\begin{mideea}{0.9}
	\begin{enumerate}
		\item \(A\vec{x} = \vec{0} \ \text{has a nontrivial solution} \ \iff \text{free variables} \iff \rule{3cm}{0.15mm}\) 

		\item A set of vectors \(\{\vec{v}_1, \vec{v}_2, \dots, \vec{v}_p \mid \vec{v}_i \in \mathbb{R}^n\}\) is linearly independent if and only if 
		\[c_1\vec{v}_1 + c_2\vec{v}_2 + \dots + c_p\vec{v}_p = \vec{0}, \ \text{with } c_i = 0, \text{for all} \ 1 \le i \le p\] is the only choice of scalars 
		\(c_i\) satisfying the vector equation. Equivalently, a set of vectors is linearly dependent if and only if there exists a vector in the set that could be 
		written as a linear combination of the other vectors.
	\end{enumerate}
\end{mideea}

\subsection{Possible/Impossible Questions}
\begin{newquestion}
	Consider the following statements. If a statement is not possible in any scenario, mark it as \textit{impossible}. Otherwise, mark it as \textit{possible}.
	\begin{table}[!h]
		\centering
		\begin{tabular}{c@{\hspace{2.5em}}cp{0.72\textwidth}}
			possible & impossible & \\
			\specialrule{\lightrulewidth}{0.25em}{1.25em}
			\bubble & \bubble & \(\{\vec{v}_1, \vec{v}_2, \vec{v}_3\}\) forms a linear independent set of vectors from \(\mathbb{R}^n\) but
						      \(\{\vec{v}_1, \vec{v}_1 + \vec{v}_2, \vec{v}_3\}\) forms a linearly dependent set. \\[0.75em]
			\bubble & \bubble & Two vectors \(\vec{v}_1, \vec{v}_2 \in \mathbb{R}^n\) form a linearly dependent set but one of the vectors
						      is not in the span of the other vector.\\[0.75em]
			\bubble & \bubble & A matrix \(A \in \mathbb{R}^{2\times3}\) where the columns of \(A\) form a linearly independent set.\\[0.75em]
			\bubble & \bubble & A matrix \(A \in \mathbb{R}^{m \times n}\) whose columns form a linearly independent set, and there exists an infinite 
						      number of solutions to the consistent system \(A\vec{x} = \vec{b}\) for some \(\vec{b} \in \mathbb{R}^m\).\\[0.75em]
			\bubble & \bubble & A linearly independent set \(\{\vec{v}_1, \vec{v}_2, \vec{v}_3\}\) where \(\{\vec{v}_1, \vec{v}_2\}\) forms a linearly
						      dependent set.\\[0.75em]
			\bottomrule
		\end{tabular}
	\end{table}
\end{newquestion}

\subsection{Short Answers Questions}
\begin{newquestion}
	For each part, provide your work (or/)and answer. 
	\begin{enumerate}
		\item Find the value(s) of \(h\) where the set of vectors \(\left\{\begin{bmatrix}-2 \\ 6 \\ 1\end{bmatrix}, \begin{bmatrix}1 \\ h \\ 2\end{bmatrix}, \begin{bmatrix}7 \\ -3 \\ 4\end{bmatrix}\right\}\) becomes linearly dependent. 
			After finding \(h\), what does this set span? (\textit{Choose only 1.})
				\[
				\begin{array}{c@{\hspace{7em}}c@{\hspace{7em}}c}
					\bubble \ \text{A plane in} \ \mathbb{R}^3 & \bubble \ \mathbb{R}^2 & \bubble \ \text{All of} \ \mathbb{R}^3
				\end{array}
				\]
		\item If, for every \(\vec{b} \in \mathbb{R}^m\), the equation \(A\vec{x} = \vec{b}\) has at most one solution, do the columns of \(A\) necessarily
			form a linearly independent set? (Note that \(A \in \mathbb{R}^{m \times n}\).)
		\item Given the matrix \(\begin{bmatrix}1 & 0 & -2 \\ 1 & 1 & -2 \\ 3 & 0 & -6\end{bmatrix}\), find a linear dependence relationship amongst the 
			column vectors, and use that to form the vector equation that shows that the columns do not form a linearly independent set.
	\end{enumerate}
\end{newquestion}

\newpage

\fancyhead[C]{\textit{Solutions}}

\textbf{Note}: The empty line under \textbf{Main Idea \#1} should be \textit{linear dependence}.

\section{Solutions to Possible/Impossible Questions}
	\begin{table}[!h]
		\centering
		\begin{tabular}{c@{\hspace{2.5em}}cp{0.72\textwidth}}
			possible & impossible & \\
			\specialrule{\lightrulewidth}{0.25em}{1.25em}
			\bubble & \dbubble & \(\{\vec{v}_1, \vec{v}_2, \vec{v}_3\}\) forms a linear independent set of vectors from \(\mathbb{R}^n\) but
						      \(\{\vec{v}_1, \vec{v}_1 + \vec{v}_2, \vec{v}_3\}\) forms a linearly dependent set. \\[0.75em]
			\dbubble & \bubble & Two vectors \(\vec{v}_1, \vec{v}_2 \in \mathbb{R}^n\) form a linearly dependent set but one of the vectors
						      is not in the span of the other vector.\\[0.75em]
			\bubble & \dbubble & A matrix \(A \in \mathbb{R}^{2\times3}\) where the columns of \(A\) form a linearly independent set.\\[0.75em]
			\bubble & \dbubble & A matrix \(A \in \mathbb{R}^{m \times n}\) whose columns form a linearly independent set, and there exists an infinite 
						      number of solutions to the consistent system \(A\vec{x} = \vec{b}\) for some \(\vec{b} \in \mathbb{R}^m\).\\[0.75em]
			\bubble & \dbubble & A linearly independent set \(\{\vec{v}_1, \vec{v}_2, \vec{v}_3\}\) where \(\{\vec{v}_1, \vec{v}_2\}\) forms a linearly
						      dependent set.\\[0.75em]
			\bottomrule
		\end{tabular}
	\end{table}


\section{Solutions to Short Answers Questions}
\begin{enumerate}
	\item \(h = 3\). To see this, row reduce the following matrix to get 
		\[\begin{bmatrix}-2 & 1 & 7 \\ 6 & h & -3 \\ 1 & 2 & 4\end{bmatrix} \sim \cdots \sim \begin{bmatrix}1 & 0 & -2 \\ 0 & 1 & 3 \\ 0 & 0 & 9-3h\end{bmatrix}.\]
		Since the matrix has only 2 pivots, it spans a plane in \(\mathbb{R}^3\).
	\item Yes. Consider the case when \(\vec{b} = \vec{0}\). That means that \(A\vec{x} = \vec{0}\) only has the trivial solution as a solution. We immediately 
		then see that the columns of \(A\) must form a linearly independent set. (If it's still hard to see, write out the matrix-vector product as a linear 
		combination of the columns of \(A\),)
	\item An example could be \[2\begin{bmatrix}1 \\ 1 \\ 3\end{bmatrix} + 0\begin{bmatrix}0 \\ 1 \\ 0\end{bmatrix} + 1\begin{bmatrix}-2 \\ -2 \\ -6\end{bmatrix} = \vec{0}.\]
\end{enumerate}


\end{document}